\chapter{Data and algorithm}\label{ch:algorithm}

\section{What has to be kept, and how much that is}\label{sec:sparsity}

Two things have to be kept: the current board, $L^2$ heights, and the record of every
avalanche, four integers each. Everything else --- the toppling matrix, the
recurrent-configuration count, the fitted exponents --- is either derived on the fly or
produced afterwards from the record.

The toppling matrix is the interesting case: it is $N \times N$ with $N = L^2$, but
\eqref{eq:delta} gives each row only five non-zero entries, so of the $N^2$ places only
about $5N$ carry information.

\begin{table}[htbp]
\centering
\small
\caption{Size of $\D$ against the number of entries it actually uses.}
\label{tab:sparsity}
\begin{tabular}{lrrr}
\toprule
\textbf{Board} & \textbf{Sites $N = L^2$} & \textbf{Places in $\D$: $N^2$}
  & \textbf{Non-zero} \\
\midrule
$L = 3$    & $9$           & $81$            & $37$ \\
$L = 50$   & $2\,500$      & $6.25\times10^{6}$ & $1.25\times10^{4}$ \\
$L = 1000$ & $10^{6}$      & $10^{12}$       & $5\times10^{6}$ \\
\bottomrule
\end{tabular}
\end{table}

At $L = 1000$ the table has a trillion places and five million useful ones. Storing it
would be absurd, and it is unnecessary, because $\D$ is not data: it is a \emph{function}
of the pair $(\bm{x},\bm{y})$, given in closed form by \eqref{eq:delta}. Looking up one
entry means asking whether $\bm{x} = \bm{y}$ and, if not, whether the two sites are
neighbours --- a constant number of comparisons, so no table is stored and the function
is evaluated where it is needed. The five entries of row $\bm{x}$ are then exactly
the diagonal subtraction and the four neighbour increments of \texttt{topple\_once}, and
the two representations are checked against each other so that they cannot drift apart.
The dense matrix is formed only for small boards and algebraic checks such as
$\det(-\D)$ (Appendix~\ref{app:matrix}).

\section{Storage and the main loop}\label{sec:storage}

The board is a one-dimensional array of length $L^2$ flattened in row-major order,
$\mathrm{idx}(\bm{x}) = iL + j$. One index rather than two is a correctness argument
rather than a micro-optimisation: toppling touches the four neighbours constantly, and
adding one to a single index is easier to verify than changing a pair of indices,
especially in the boundary test where an unsigned underflow would silently index the
wrong site. Each avalanche appends four integers in order with no summary column,
because the later binning, fitting and correlation analysis all need the raw rows.

\begin{lstlisting}[caption={Main loop: burn-in, sampling, writing the four columns.}]
SandPile<BtwMatrix> pile(L, L, seed);
for (std::uint64_t t = 0; t < 2 * L * L; ++t) pile.drop_sand();   // burn-in
pile.clear_history();
std::ofstream out{"avalanches.csv"};
out << "t,s,area,loss,length\n";
for (std::uint64_t t = 1; t <= drops; ++t) {
  pile.drop_sand();
  const AvalancheStats& st = pile.history().back();
  out << t << ',' << st.topples << ',' << st.area << ','
      << st.loss << ',' << st.length << '\n';
}
\end{lstlisting}

Randomness enters in one place only, the uniformly drawn drop site, so the avalanche is
deterministic and a fixed seed reproduces a run exactly. That matters for
Chapter~\ref{ch:hexagonal}, where the difference between two lattices must not be a
difference between two random streams. The cost of a drop is $O(1)$ plus $O(s)$ for the
avalanche, so runtime is proportional to the number of topplings, which is the quantity
being measured.

\section{Burn-in}\label{sec:burnin}

The early part of a run has to be discarded. The board starts empty, so large avalanches
are rare while it fills, and including that stage depresses the tail of the distribution
and biases any exponent fitted from it upward. The run is therefore driven for $2L^2$
drops and the history cleared before sampling begins.

Whether the board is in the stationary state is checked by quantities that need no
fitting and no choice of cut-off: the mean height against the exact value $17/8 = 2.125$
of \eqref{eq:meanh}, and the mean number of grains lost per drop against $1$, which is
the input--output balance of §\ref{sec:noneq}. Both are sensitive to the two errors that
would invalidate everything downstream --- a wrong boundary test and a wrong threshold
--- which is why they are checked before any avalanche statistic is recorded.

\begin{pending}
  \pendingtag The runs whose statistics are reported in Chapter~\ref{ch:stationary} have
  not yet been carried out. What is needed is a set of runs at three or more board sizes
  with a common seed policy and a common burn-in, each long enough that the usable tail
  of the toppling distribution holds on the order of $10^3$ events; the record is then
  sufficient for both the exponent fitting of §\ref{sec:distributions} and the
  conditional means of §\ref{sec:correlations}.
\end{pending}
